\begin{lemma}
\label{lemma-orbits}
Let $k$ be a field. Let $n \geq 2$. For $1 \leq i, j \leq n$ with
$i \not = j$ and $d \geq 0$ denote $T_{i, j, d}$ the automorphism
of $\mathbf{A}^n_k$ given in coordinates by
$$
(x_1, \ldots, x_n) \longmapsto
(x_1, \ldots, x_{i - 1}, x_i + x_j^d, x_{i + 1}, \ldots, x_n)
$$
Let $W \subset \mathbf{A}^n_k$ be a nonempty open subscheme
such that $T_{i, j, d}(W) = W$ for all $i, j, d$ as above.
Then either $W = \mathbf{A}^n_k$ or the characteristic of $k$
is $p > 0$ and $\mathbf{A}^n_k \setminus W$ is a finite set
of closed points whose coordinates are algebraic over $\mathbf{F}_p$.
\end{lemma}

\begin{proof}
We may replace $k$ by any extension field in order to prove this.
Let $Z$ be an irreducible component of $\mathbf{A}^n_k \setminus W$.
Assume $\dim(Z) \geq 1$, to get a contradiction.
Then there exists an extension field $k'/k$ and a $k'$-valued
point $\xi = (\xi_1, \ldots, \xi_n) \in (k')^n$ of
$Z_{k'} \subset \mathbf{A}^n_{k'}$
such that at least one of $x_1, \ldots, x_n$ is transcendental over the
prime field. Claim: the orbit of $\xi$ under the group generated by
the transformations $T_{i, j, d}$ is Zariski
dense in $\mathbf{A}^n_{k'}$. The claim will give the desired contradiction.

\medskip\noindent
If the characteristic of $k'$ is zero, then already the operators
$T_{i, j, 0}$ will be enough since these transform $\xi$ into
the points
$$
(\xi_1 + a_1, \ldots, \xi_n + a_n)
$$
for arbitrary $(a_1, \ldots, a_n) \in \mathbf{Z}_{\geq 0}^n$.
If the characteristic is $p > 0$, we may assume after renumbering
that $\xi_n$ is transcendental over $\mathbf{F}_p$. By
successively applying the operators $T_{i, n, d}$ for
$i < n$ we see the orbit of $\xi$ contains the elements
$$
(\xi_1 + P_1(\xi_n), \ldots, \xi_{n - 1} + P_{n - 1}(\xi_n), \xi_n)
$$
for arbitrary $(P_1, \ldots, P_{n - 1}) \in \mathbf{F}_p[t]$.
Thus the Zariski closure of the orbit contains the coordinate
hyperplane $x_n = \xi_n$. Repeating the argument with a different
coordinate, we conclude that the Zariski closure contains
$x_i = \xi_i + P(\xi_n)$ for any $P \in \mathbf{F}_p[t]$
such that $\xi_i + P(\xi_n)$ is transcendental over $\mathbf{F}_p$.
Since there are infinitely many such $P$ the claim follows.

\medskip\noindent
Of course the argument in the preceding paragraph also applies
if $Z = \{z\}$ has dimension $0$ and the coordinates of $z$
in $\kappa(z)$ are not algebraic over $\mathbf{F}_p$. The lemma follows.
\end{proof}